\subsection{同じ近似分数を,整数の更新から見直す}\label{節：ペル方程式と近似分数}
  \sectionref*{節：ニュートン法型}で,$a_1=1$から$\sqrt2$を近似すると,
  \[
    a_2=\frac32,\qquad a_3=\frac{17}{12},\qquad a_4=\frac{577}{408}
  \]
  が現れた.ここでは分子と分母を別々に追い,
  連立型・一次分数型・ニュートン法が同じ分数に結び付くことを確かめる.
  微積分を使わず,二次式の計算と極限で追える内容である.

  \bookstep{分子と分母に残る関係を見る}
  $3^2-2\cdot2^2=17^2-2\cdot12^2=577^2-2\cdot408^2=1$である.
  $x^2-dy^2=1$の整数解を求める問題を\textbf{ペル方程式}という.
  ここでは$d=2$の具体例を扱い,一般の解の存在や全解の分類には進まない.
  次の数列を作ろう.
  \[
    x_k+y_k\sqrt2=(3+2\sqrt2)^k\quad(k\ge1)
  \]
  両辺を$3+2\sqrt2$倍して有理部分と$\sqrt2$の係数を比べると,
  \[
    x_{k+1}=3x_k+4y_k,\qquad y_{k+1}=2x_k+3y_k
  \]
  となる.$x_1=3$, $y_1=2$なので,両方とも正の整数である.
  共役な式との積から
  \[
    x_k^2-2y_k^2=(3+2\sqrt2)^k(3-2\sqrt2)^k=1
  \]
  が保たれる.整数の式を保つ更新は,
  \bookproblemref{practice-bridge}{9}{発展入門問9}の見方とも共通する.

  \bookstep{比を取ると,一次分数型になる}
  $r_k=\frac{x_k}{y_k}$とすると,連立の更新から
  \[
    r_{k+1}=\frac{3r_k+4}{2r_k+3}
  \]
  である.分母は正であり,不動点を求めると$\pm\sqrt2$になる.
  一次分数型で用いた差の比は,
  \[
    \frac{r_{k+1}-\sqrt2}{r_{k+1}+\sqrt2}
    =(3-2\sqrt2)^2
      \frac{r_k-\sqrt2}{r_k+\sqrt2}
  \]
  と更新される.整数の連立更新と比の一次分数更新は,
  同じ計算を別々の量で追ったものだった.

  \bookstep{ニュートン法では,添字を倍に進める}
  $(x_k+y_k\sqrt2)^2$の展開から,
  \[
    x_{2k}=x_k^2+2y_k^2,\qquad y_{2k}=2x_ky_k
  \]
  なので,
  \[
    r_{2k}=\frac12\left(r_k+\frac2{r_k}\right).
  \]
  これはニュートン法型の更新である.従って本編の$a_n$は
  \[
    a_n=r_{2^{n-2}}\qquad(n\ge2)
  \]
  に当たる.本編の$3/2,17/12,577/408$は$r_1,r_2,r_4$であり,
  この整数更新の連続する三項ではない.
  \BookFigPellDoubling
  $x_k^2-2y_k^2=1$を因数分解すれば,
  \[
    r_k-\sqrt2=\frac1{y_k^2(r_k+\sqrt2)}>0.
  \]
  $y_{k+1}>3y_k$なので$y_k\to\infty$であり,$r_k\to\sqrt2$となる.
  ニュートン法の速さは,この列で添字を倍に進めることからも読み取れる.

  \bookstep{連分数の有限段階にも,同じ分数が現れる}
  $\sqrt2=1+\frac1{1+\sqrt2}$から,分母に$2$が繰り返される表示を考える.
  無限の式を先に仮定せず,有限段階の近似を
  \[
    C_0=1,\qquad C_{j+1}=1+\frac1{1+C_j}
    \quad(j\ge0)
  \]
  で定める.$C_1=3/2$, $C_2=7/5$, $C_3=17/12$である.
  $C_j=P_j/Q_j$とすると,$P_0=Q_0=1$から
  \[
    P_{j+1}=P_j+2Q_j,\qquad Q_{j+1}=P_j+Q_j
  \]
  が得られ,$P_j+Q_j\sqrt2=(1+\sqrt2)^{j+1}$となる.
  $(1+\sqrt2)^2=3+2\sqrt2$だから,
  \[
    C_{2k-1}=r_k\qquad(k\ge1)
  \]
  である.この有限段階の値を\textbf{連分数の近似分数}と呼ぶ.
  偶数番号では$P_j^2-2Q_j^2=-1$,奇数番号では$1$となり,
  同じ誤差の因数分解で$C_j\to\sqrt2$も確認できる.

  \begin{bookpoint}
    \textbf{三つの解法を,同じ対象へ戻す}\par
    分子・分母を追えば連立型,比を追えば一次分数型,
    添字を倍に進めれば平方根のニュートン法になる.
    形式的な変形の違いは,何を追い,何回分の更新を見るかの違いでもある.
  \end{bookpoint}
  保存される式と誤差を\bookproblemref{practice-connections}{6}{接続演習問6}で確かめよう.
  ペル方程式と連分数の補足資料は,参考文献のKeith Conradの二つの解説へ.
